\section{Einleitung}
\subsection{Wozu Git?}

\begin{frame}{Wozu Git?}
    \uncover<0->{Lasst uns zusammen ein Programm schreiben!}
    \uncover<2->{\textit{(Evolution eines Hackers)}}

    \begin{enumerate}
        \item<3-> Wir arbeiten im Etherpad!
        \item<4-> Wir arbeiten auf \texttt{rzssh1}!
        \item<5-> Du schickst mir dann die neue Version per Mail!
        \item<6-> Du schickst mir dann ein Diff per Mail!
        \item<7-> Wir nehmen Versionskontrolle!
    \end{enumerate}

    \uncover<8> {
        $\rightarrow$ gilt auch für Hausaufgaben in \LaTeX
    }

\iffinal
    \vspace{0.2cm}

    \makebox[\textwidth][c]{
        \begin{tikzpicture}[x=1cm,y=1cm]
            \uncover<3->{
                \node[anchor=south west,inner sep=0,opacity=0.2] (human-01) at (0,0) {\includegraphics[scale=0.5]{gfx/human-01.png}};
            }
            \uncover<4->{
                \node[anchor=south west,inner sep=0,opacity=0.2] (human-02) at (1.2,0) {\includegraphics[scale=0.5]{gfx/human-02.png}};
            }
            \uncover<5->{
                \node[anchor=south west,inner sep=0,opacity=0.2] (human-03) at (2.1,0) {\includegraphics[scale=0.5]{gfx/human-03.png}};
            }
            \uncover<6->{
                \node[anchor=south west,inner sep=0,opacity=0.2] (human-04) at (2.9,0) {\includegraphics[scale=0.5]{gfx/human-04.png}};
            }
            \uncover<7->{
                \node[anchor=south west,inner sep=0,opacity=0.2] (human-05) at (4.0,0) {\includegraphics[scale=0.5]{gfx/human-05.png}};
            }
        \end{tikzpicture}
    }
\fi
\end{frame}

\begin{frame}{Was ist Versionierung?}
    \begin{itemize}
        \item \emph{Snapshots} einzelner Dateiversionen speichern
        \item History aufbewahren und wiederherstellbar machen
        \item verschiedene Versionen zusammenführen
        \item \emph{(optional)} Synchronisation mit entfernten Kopien (Kollaboration)
    \end{itemize}
    $\rightarrow$ mehr als nur ein Backup
\end{frame}

\subsection{Warum gerade git?}
\begin{frame}{Warum gerade git?}
    Vorteile:
    \begin{itemize}
        \item verteilt (serverunabhängig, jeder \emph{clone} ist eigenständig/vollständig)
        \item schnell (lokal, Implementation in C)
        \item optimal für Quelltext
        \item FOSS
    \end{itemize}

    Nachteile:
    \begin{itemize}
        \item wenig geeignet für Binärdateien
        \item gesamte History kann groß werden
    \end{itemize}
\end{frame}
